\DocumentMetadata{
  pdfversion=2.0,
  pdfstandard=ua-2,
  lang=en-US,
  tagging=on,
  tagging-setup={math/setup=mathml-SE,tabsorder=structure}
}
\documentclass[11pt,letterpaper]{article}
\usepackage[margin=0.68in,includefoot]{geometry}
\usepackage{newcomputermodern}
\usepackage{amsmath,amscd,mathtools}
\usepackage{tikz,adjustbox}
\providecommand{\square}{\mdlgwhtsquare}
\usepackage[hidelinks]{hyperref}
\hypersetup{
  pdftitle={Combinatorics PhD Exam, May 2021},
  pdfauthor={University of Florida Department of Mathematics},
  pdfsubject={Graduate mathematics examination},
  pdfdisplaydoctitle=true,
  bookmarksopen=true
}
\setcounter{secnumdepth}{0}
\setlength{\parindent}{0pt}
\setlength{\parskip}{0.28em}
\setlength{\emergencystretch}{3em}
\setlength{\leftmargini}{2.15em}
\setlength{\leftmarginii}{2.65em}
\setlength{\leftmarginiii}{3em}
\pagestyle{empty}
\raggedbottom
\allowdisplaybreaks
\NewTaggingSocketPlug{block/list/label}{examproblem}{%
  \tagstructbegin{tag=Lbl}\tagstructend%
  \tagstructbegin{tag=\LBody}%
  \tagstructbegin{tag=\ProblemHeadingTag,title-o={\ProblemHeadingText}}%
  \tagmcbegin{tag=\ProblemHeadingTag,actualtext={\ProblemHeadingText}}%
  #2%
  \tagmcend%
  \tagstructend%
}
\newenvironment{examproblems}[2]{%
  \def\ProblemHeadingTag{#2}%
  \AssignTaggingSocketPlug{block/list/label}{examproblem}%
  \begin{enumerate}%
  \setcounter{enumi}{#1}%
  \renewcommand{\labelenumi}{\textbf{\theenumi.}}%
  \setlength{\topsep}{0.35em}%
  \setlength{\partopsep}{0pt}%
  \setlength{\itemsep}{0.48em}%
  \setlength{\parsep}{0.18em}%
}{\end{enumerate}}
\newenvironment{examsubparts}{%
  \AssignTaggingSocketPlug{block/list/label}{default}%
  \begin{enumerate}%
  \setlength{\topsep}{0.18em}%
  \setlength{\partopsep}{0pt}%
  \setlength{\itemsep}{0.16em}%
  \setlength{\parsep}{0.1em}%
}{\end{enumerate}}
\newenvironment{examinstructions}{%
  \par\begingroup\itshape
}{%
  \par\endgroup\vspace{0.18em}
}
\makeatletter
\renewcommand\subsection{\@startsection{subsection}{2}{0pt}%
  {-1.25ex plus -.25ex minus -.1ex}%
  {0.55ex plus .1ex}%
  {\normalfont\large\bfseries}}
\renewcommand{\@maketitle}{%
  \newpage\null\vskip -1em%
  {\footnotesize\bfseries
    \href{https://www.ufl.edu/}{University of Florida}, \href{https://math.ufl.edu/}{Department of Mathematics}\par}%
  \vskip -0.5em%
  \tagstructbegin{tag=H1,title={Combinatorics PhD Exam, May 2021}}%
  \tagpdfparaOff%
  \tagmcbegin{tag=H1}%
    {\LARGE\bfseries\href{https://gma.math.ufl.edu/exams/combinatorics-phd/}{Combinatorics PhD Exam}, May 2021\par}%
  \tagmcend%
  \tagpdfparaOn%
  \tagstructend%
  \vskip 0.25em%
  \tagmcbegin{artifact}%
    \hrule height 1pt%
  \tagmcend%
  \vskip 1em%
}
\makeatother
\title{Combinatorics PhD Exam, May 2021}
\author{}
\date{}
\begin{document}
\maketitle
\thispagestyle{empty}
\begin{examproblems}{0}{H2}
\def\ProblemHeadingText{Problem 1.}
\item
Let \(a_n\) be the number of ways to select a permutation of length~\(n\), and then color each of its cycles red, blue, or green. Find an explicit formula for \(a_n\).
\def\ProblemHeadingText{Problem 2.}
\item
Let~\(A\) be the graph obtained from \(K_n\) by deleting an edge. Find a closed formula (so no summation signs) for the number of spanning trees of~\(A\).
\def\ProblemHeadingText{Problem 3.}
\item
Prove bijectively that the number of graphs with vertex set \([n]\) for which all vertices have even degree is \(2^{\binom{n-1}{2}}\).
\def\ProblemHeadingText{Problem 4.}
\item
Prove that for every~\(n\), there is a tournament~\(T\) on~\(n\) vertices with at least \(n!/2^{n-1}\) Hamiltonian paths.
\def\ProblemHeadingText{Problem 5.}
\item
Let \(\mathbb{P}=\{1,2,\ldots\}\) and \((\mathbb{P},\mid)\) be the division poset (so \(m\mid n\) if \(n/m\in\mathbb{P}\)). Find a formula for the Möbius function~\(\mu\) and prove that it is correct. (Hint: decompose intervals in \((\mathbb{P},\mid)\) as products of chains.)
\def\ProblemHeadingText{Problem 6.}
\item
Let~\(\mathcal{A}\) be the class of compositions into odd parts, where the size of \(\alpha=(a_1,\ldots,a_n)\in\mathcal{A}\) is \(|\alpha|=a_1+\cdots+a_n\). Find the associated ordinary generating function \(A(x)=\sum_{n\geq 0}a_nx^n\) and compute its growth rate \(\limsup a_n^{1/n}\).
\def\ProblemHeadingText{Problem 7.}
\item
Let~\(\mathcal{P}\) be the class of integer partitions. Give a bijective proof of the Cauchy identity
\[
\prod_{i,j\geq 1}(1-x_i y_j)^{-1}=\sum_{\lambda\in\mathcal{P}}s_\lambda(x_1,x_2,\ldots)s_\lambda(y_1,y_2,\ldots)
\]
 where \(s_\lambda\) is the Schur function of shape~\(\lambda\).
\def\ProblemHeadingText{Problem 8.}
\item
Let~\(\mathcal{A}\) and~\(\mathcal{B}\) be labeled combinatorial classes with size functions \(|\cdot|_{\mathcal{A}}\) and \(|\cdot|_{\mathcal{B}}\), respectively. Define an admissible construction~\(\star\) so that \(\mathcal{C}=\mathcal{A}\star\mathcal{B}\) satisfies the exponential generating function identity \(C_E(x)=A_E(x)\cdot B_E(x)\).
\end{examproblems}
\end{document}
